% !TEX program = pdflatex
\documentclass[aspectratio=169,10pt]{beamer}

\usepackage[T1]{fontenc}
\usepackage[utf8]{inputenc}
\usepackage[spanish]{babel}
\usepackage{booktabs}
\usepackage{tabularx}
\usepackage{graphicx}
\usepackage{hyperref}
\usepackage{xcolor}
\usepackage{tikz}
\usepackage{enumitem}
\setlist{leftmargin=*,itemsep=0.22em,topsep=0.15em}

\usetheme{Madrid}
\usecolortheme{default}
\setbeamertemplate{navigation symbols}{}
\setbeamertemplate{footline}{%
  \hfill{\scriptsize\color{black!60}\insertsectionhead\ \textbar\ \insertframenumber/\inserttotalframenumber}\hspace{1em}\vspace{0.35em}%
}

\definecolor{ieblue}{RGB}{0,51,102}
\definecolor{good}{RGB}{34,120,70}
\definecolor{bad}{RGB}{160,40,40}
\definecolor{soft}{RGB}{245,247,250}
\definecolor{accent}{RGB}{0,120,160}
\definecolor{esa}{RGB}{0,70,130}
\definecolor{phb}{RGB}{180,90,20}
\setbeamercolor{structure}{fg=ieblue}
\setbeamercolor{frametitle}{fg=ieblue,bg=soft}
\setbeamercolor{block title}{bg=ieblue,fg=white}
\setbeamercolor{block body}{bg=soft}
\setbeamerfont{frametitle}{size=\large,series=\bfseries}

% Shared TikZ helpers --- layouts seguros (sin solapes)
\usetikzlibrary{arrows.meta,positioning,shapes.geometric,calc}

\tikzset{
  box/.style={
    draw=ieblue, rounded corners=3pt, align=center,
    fill=soft, font=\scriptsize, inner sep=5pt,
    minimum height=8.5mm, text depth=0.3em
  },
  boxdark/.style={
    draw=ieblue, rounded corners=3pt, align=center,
    fill=ieblue, text=white, font=\scriptsize\bfseries, inner sep=5pt,
    minimum height=8.5mm, text depth=0.3em
  },
  arr/.style={-{Latex}, thick, ieblue},
  pill/.style={
    draw=ieblue, rounded corners=8pt, fill=white,
    font=\scriptsize\bfseries, inner sep=4pt
  }
}

% Pipeline horizontal con anchos fijos y hueco para flechas
\newcommand{\PipelineFive}[5]{%
\begin{center}
\begin{tikzpicture}
  \node[box, text width=1.85cm] (a) {#1};
  \node[box, text width=1.85cm, right=0.55cm of a] (b) {#2};
  \node[box, text width=1.85cm, right=0.55cm of b] (c) {#3};
  \node[box, text width=1.85cm, right=0.55cm of c] (d) {#4};
  \node[boxdark, text width=1.85cm, right=0.55cm of d] (e) {#5};
  \draw[arr] (a) -- (b);
  \draw[arr] (b) -- (c);
  \draw[arr] (c) -- (d);
  \draw[arr] (d) -- (e);
\end{tikzpicture}
\end{center}%
}

% Encaje Pedro / Mario con minipages (evita solapes TikZ)
\newcommand{\FitDual}[4]{%
\noindent
\begin{minipage}[t]{0.48\textwidth}
\centering
\begin{tikzpicture}\node[circle, draw=ieblue, fill=blue!12, minimum size=1cm, font=\small\bfseries]{Pedro};\end{tikzpicture}\\[0.25em]
\begin{beamercolorbox}[wd=\textwidth,rounded=true,sep=4pt]{block body}
\scriptsize #1\\[0.35em]\textbf{Afinidad: #3/10}
\end{beamercolorbox}
\end{minipage}\hfill
\begin{minipage}[t]{0.48\textwidth}
\centering
\begin{tikzpicture}\node[circle, draw=orange!70!black, fill=orange!18, minimum size=1cm, font=\small\bfseries]{Mario};\end{tikzpicture}\\[0.25em]
\begin{beamercolorbox}[wd=\textwidth,rounded=true,sep=4pt]{block body}
\scriptsize #2\\[0.35em]\textbf{Afinidad: #4/10}
\end{beamercolorbox}
\end{minipage}%
}

% Pros / contras con minipages
\newcommand{\ProCon}[2]{%
\noindent
\begin{minipage}[t]{0.48\textwidth}
\begin{beamercolorbox}[wd=\textwidth,rounded=true,sep=5pt]{block body}
\textcolor{good}{\textbf{Lo bueno}}\\[0.3em]
\scriptsize #1
\end{beamercolorbox}
\end{minipage}\hfill
\begin{minipage}[t]{0.48\textwidth}
\begin{beamercolorbox}[wd=\textwidth,rounded=true,sep=5pt]{block body}
\textcolor{bad}{\textbf{Lo malo}}\\[0.3em]
\scriptsize #2
\end{beamercolorbox}
\end{minipage}%
}

% Temario vs extensión
\newcommand{\MethodsSplit}[2]{%
\noindent
\begin{minipage}[t]{0.48\textwidth}
\begin{beamercolorbox}[wd=\textwidth,rounded=true,sep=5pt]{block body}
\textbf{Del curso}\\[0.25em]\scriptsize #1
\end{beamercolorbox}
\end{minipage}\hfill
\begin{minipage}[t]{0.48\textwidth}
\begin{beamercolorbox}[wd=\textwidth,rounded=true,sep=5pt]{block title}
\textbf{Más allá del temario}\\[0.15em]
\end{beamercolorbox}
\begin{beamercolorbox}[wd=\textwidth,rounded=true,sep=5pt]{block body}
\scriptsize #2
\end{beamercolorbox}
\end{minipage}%
}

% Beamercolorbox (no TikZ): evita solapes cuando el subtítulo es largo
\newcommand{\Verdict}[2]{%
\begin{center}
\begin{beamercolorbox}[wd=0.92\textwidth,rounded=true,sep=10pt,center]{block title}
\textbf{\large #1}\\[0.45em]
{\normalsize\normalfont #2}
\end{beamercolorbox}
\end{center}%
}

\newcommand{\ScoreBar}[3]{%
\noindent\begin{tikzpicture}
  \node[anchor=west, font=\scriptsize, text width=2.3cm] at (0,0) {#1};
  \fill[black!10] (2.4, -0.12) rectangle (6.6, 0.22);
  \fill[#3] (2.4, -0.12) rectangle ({2.4+#2*0.42}, 0.22);
  \node[font=\scriptsize\bfseries, anchor=west] at (6.75,0.05) {#2};
\end{tikzpicture}\par\vspace{0.1em}%
}

% Cabecera de opción sin solapes
\newcommand{\OptionHeader}[3]{%
\begin{center}
\begin{tikzpicture}
  \node[boxdark, text width=4.0cm] (a) {#1};
  \node[pill, right=0.5cm of a] (b) {#2};
\end{tikzpicture}
\end{center}
\vspace{0.35em}
\begin{block}{En una frase}
#3
\end{block}%
}


\title{ESA vs PHBench}
\subtitle{Decision deck --- shortlist final Pedro $\times$ Mario}
\author{Deep research $\cdot$ Octubre 2026}
\date{11 octubre 2026}
\institute{BAM $\cdot$ Machine Learning \& Communication Skills $\cdot$ IE University}

\begin{document}

\begin{frame}{Portada}
  \titlepage
\end{frame}

\begin{frame}{Para qué sirve este deck}
  \begin{columns}[T]
    \begin{column}{0.55\textwidth}
      \begin{itemize}
        \item Ya acotamos a \textbf{dos} opciones finales.
        \item Aquí comparamos en profundidad, sin jerga innecesaria.
        \item Objetivo: elegir \textbf{1 principal + 1 reserva} esta noche.
        \item Si una no tiene datos mañana, no puede ser principal.
      \end{itemize}
    \end{column}
    \begin{column}{0.42\textwidth}
      \begin{tikzpicture}
        \node[boxdark, text width=4.4cm] (a) {1. Entender cada una};
        \node[box, text width=4.4cm, below=0.28cm of a] (b) {2. Comparar cara a cara};
        \node[box, text width=4.4cm, below=0.28cm of b] (c) {3. Escenarios de decisión};
        \node[boxdark, text width=4.4cm, below=0.28cm of c] (d) {4. Principal + reserva};
        \draw[arr] (a)--(b); \draw[arr] (b)--(c); \draw[arr] (c)--(d);
      \end{tikzpicture}
    \end{column}
  \end{columns}
\end{frame}

\begin{frame}{Las dos finalistas (en una frase)}
  \begin{columns}[T]
    \begin{column}{0.48\textwidth}
      \begin{beamercolorbox}[wd=\textwidth,rounded=true,sep=6pt]{block title}
        \textbf{A · ESA Collision Avoidance}
      \end{beamercolorbox}
      \begin{beamercolorbox}[wd=\textwidth,rounded=true,sep=6pt]{block body}
        \scriptsize Con avisos tempranos de ``casi choque'' en órbita, anticipar si el riesgo final será grave --- con datos oficiales ESA.
      \end{beamercolorbox}
    \end{column}
    \begin{column}{0.48\textwidth}
      \begin{beamercolorbox}[wd=\textwidth,rounded=true,sep=6pt]{block title}
        \textbf{B · PHBench (Series A)}
      \end{beamercolorbox}
      \begin{beamercolorbox}[wd=\textwidth,rounded=true,sep=6pt]{block body}
        \scriptsize Con señales públicas del lanzamiento en Product Hunt, ordenar qué startups merecen atención por posible Series A.
      \end{beamercolorbox}
    \end{column}
  \end{columns}
  \vspace{0.6em}
  \begin{center}
  \begin{tikzpicture}
    \node[boxdark, text width=3.8cm, fill=esa] {Espacio / deeptech};
    \node[font=\Large\bfseries] at (5.2,0) {vs};
    \node[boxdark, text width=3.8cm, fill=phb, xshift=7.2cm] {VC / ranking};
  \end{tikzpicture}
  \end{center}
\end{frame}

\begin{frame}{Filtro del curso (igual para las dos)}
  \begin{center}
  \begin{tikzpicture}
    \node[boxdark, text width=2.5cm] (d) {Datos reales\\manejables};
    \node[boxdark, text width=2.5cm, right=0.45cm of d] (m) {Métodos del\\temario (+extra)};
    \node[boxdark, text width=2.5cm, right=0.45cm of m] (a) {Ambición\\visible (25\%)};
    \node[boxdark, text width=2.5cm, right=0.45cm of a] (o) {Defendible\\en oral};
  \end{tikzpicture}
  \end{center}
  \vspace{0.45em}
  \begin{block}{Marco práctico}
    Python + scikit-learn + PyTorch \quad$\cdot$\quad $\sim$6 semanas \quad$\cdot$\quad informe + GitHub con historial real \quad$\cdot$\quad proposal $\sim$2 páginas \textbf{ahora}
  \end{block}
\end{frame}

\begin{frame}{Filtro nuestro (Pedro $\times$ Mario)}
  \FitDual{Deeptech, física, espacio, quantum. Rechaza lo genérico. Quiere impresionar. No busca el camino fácil.}{Finanzas y sistemas de recomendación. Quiere originalidad, buenos datos y comparación limpia.}{---}{---}
  \vspace{0.4em}
  \begin{block}{Vetos}
    \textbf{Mario:} nada de salud / biomedicina. \quad
    \textbf{Pedro:} nada de proyecto genérico de estudiante (crédito, house prices, bolsa típica).
  \end{block}
\end{frame}

% ============================================================
\section{Opción A --- ESA}
% ============================================================

\begin{frame}{A · Portada --- Evitar choques entre satélites}
  \OptionHeader{Finalista A}{Espacio / ESA}{%
    Usamos avisos oficiales de la Agencia Espacial Europea para anticipar si dos objetos en órbita van a acercarse de forma peligrosa.
  }
\end{frame}

\begin{frame}{A · Imagina esto (sin jerga)}
  \begin{columns}[T]
    \begin{column}{0.55\textwidth}
      \begin{itemize}
        \item El espacio cercano es una \textbf{autopista muy llena}.
        \item Satélites y basura espacial circulan a miles de km/h.
        \item Cuando dos objetos se van a cruzar cerca, llega un \textbf{aviso} (CDM).
        \item Ese aviso se actualiza durante días: el riesgo sube o baja.
        \item Alguien decide: \textbf{¿desvío el satélite o no?}
      \end{itemize}
    \end{column}
    \begin{column}{0.42\textwidth}
      \begin{tikzpicture}
        \draw[thick, ieblue] (0,0) circle (1.1);
        \node[font=\tiny] at (0,0) {Tierra};
        \draw[thick, accent] (0,0) circle (1.65);
        \fill[ieblue] (1.5,0.65) circle (0.1);
        \fill[bad] (1.85,0.4) circle (0.09);
        \node[font=\tiny, align=center] at (0.9,-1.95) {dos objetos\\se acercan};
      \end{tikzpicture}
    \end{column}
  \end{columns}
\end{frame}

\begin{frame}{A · Analogía cotidiana}
  \begin{center}
  \begin{tikzpicture}
    \node[box, text width=3.4cm] (a) {Parte del tiempo\\a las 10:00};
    \node[box, text width=3.4cm, right=0.55cm of a] (b) {¿Tormenta fuerte\\a las 18:00?};
    \node[boxdark, text width=3.4cm, right=0.55cm of b] (c) {Decidir si\\cancelar el plan};
    \draw[arr] (a)--(b)--(c);
  \end{tikzpicture}
  \end{center}
  \vspace{0.45em}
  Con el aviso \textbf{temprano} (p.\ ej.\ $\geq$2 días antes), ¿podemos estimar el riesgo \textbf{final}?\\
  Si avisamos bien y pronto: menos maniobras inútiles y menos peligros ignorados.
\end{frame}

\begin{frame}{A · ¿Qué haría el machine learning?}
  \begin{itemize}
    \item Hoy hay física + operadores humanos mirando muchos avisos.
    \item Nosotros: con la secuencia truncada, estimar el riesgo del último aviso.
    \item Éxito = detectar bien los casos \textbf{raros y peligrosos}, no acertar el 99\% fácil.
    \item Extra de ambición: \textbf{calibración}, curvas de aviso temprano, o variables físicas explícitas.
  \end{itemize}
\end{frame}

\begin{frame}{A · Glosario rápido}
  \begin{description}[style=nextline,leftmargin=1.4cm,font=\bfseries\scriptsize]
    \item[CDM] Mensaje de ``casi choque'': una fila con distancia, incertidumbre y riesgo.
    \item[Evento] Toda la secuencia de avisos del mismo acercamiento.
    \item[TCA] Momento del máximo acercamiento.
    \item[High-risk] Umbral operativo típico: $\log_{10} r \geq -6$ (riesgo $\geq 10^{-6}$).
    \item[LRP] Baseline fuerte: ``usa el último riesgo conocido $\geq$2 días antes''.
  \end{description}
\end{frame}

\begin{frame}{A · Pregunta de investigación}
  \begin{block}{Principal}
    Si solo vemos parte de la secuencia (hasta $\geq$2 días antes del TCA), ¿predecimos el riesgo final mejor que reglas simples (LRP)?
  \end{block}
  \begin{block}{Secundarias (elegid 1 como gancho)}
    \scriptsize
    (i) ¿Cómo cae la calidad al avisar cada vez antes? \quad
    (ii) ¿Las probabilidades están bien calibradas? \quad
    (iii) ¿Hay transferencia entre misiones? \quad
    (iv) ¿Ayudan variables físicas explícitas?
  \end{block}
\end{frame}

\begin{frame}{A · Datos (hechos concretos)}
  \begin{itemize}
    \item Zenodo público: \texttt{zenodo.org/records/4463683} --- \textbf{descarga inmediata} (CC BY 4.0).
    \item Train: 162\,634 filas / 13\,154 eventos \quad·\quad Test: 24\,484 / 2\,167.
    \item $\sim$103 variables por CDM; riesgo en $\log_{10}$.
    \item High-risk ($r\geq 10^{-6}$): $\sim$3{,}4\% en train; test enriquecido $\sim$6{,}9\%.
    \item Paper: arXiv:2008.03069. Winners: reglas metric-aware o LSTM siamesa --- no DL genérico.
  \end{itemize}
\end{frame}

\begin{frame}{A · Baseline que hay que batir (o explicar)}
  \begin{center}
  \begin{tikzpicture}
    \node[box, text width=3.6cm] (c) {Constante $r=-5$\\{\scriptsize $L\approx 2.50$}};
    \node[boxdark, text width=3.6cm, right=0.7cm of c] (l) {Latest Risk (LRP)\\{\scriptsize $L\approx 0.694$}};
    \node[box, text width=3.6cm, right=0.7cm of l] (w) {Solo $\sim$12/97\\equipos lo batieron};
    \draw[arr] (c)--(l); \draw[arr] (l)--(w);
  \end{tikzpicture}
  \end{center}
  \vspace{0.45em}
  \textbf{Implicación:} no basta con ``entrenar un XGBoost''. El valor está en el \textbf{ángulo} (early-warning, calibración, física, transfer).
\end{frame}

\begin{frame}{A · Qué se intentó ya (competición Kelvins 2019)}
  \begin{itemize}
    \item Reto oficial ESA (Kelvins): $\sim$97 equipos; paper Uriot et al.\ (arXiv:2008.03069).
    \item La regla ingenua \textbf{LRP} (último riesgo $\geq$2 días antes) fue mejor que el 85\% de los equipos.
    \item Ganadores no fueron ``redes profundas genéricas'':
    \begin{itemize}
      \item \textbf{1º sesc (Honda RI):} cascada de reglas sobre el riesgo a $-$2 días + recortes metric-aware.
      \item \textbf{3º Magpies:} LSTM siamesa para detectar saltos low$\to$high + ensemble.
    \end{itemize}
    \item Muchos probaron AutoML / RF / tsfresh y \textbf{sobreajustaron}.
  \end{itemize}
\end{frame}

\begin{frame}{A · Qué queda abierto (no es cosplay de leaderboard)}
  \begin{itemize}
    \item La competición optimizó una loss rara ($L = \mathrm{MSE}_{HR}/F_2$). Nosotros podemos elegir un claim científico distinto.
    \item Aún poco explorado con rigor: \textbf{curvas de lead-time}, \textbf{calibración/UQ}, leave-mission-out, ablación física limpia.
    \item Oral fuerte: ``citamos a sesc/Magpies como baseline; nuestra contribución es X''.
    \item \textbf{No} afirmar: ``vamos a ganar Kelvins'' ni ``XGBoost nuevo''.
  \end{itemize}
\end{frame}

\begin{frame}{A · Cómo lo haríamos (6 semanas)}
  \PipelineFive{EDA +\\baselines}{Features\\físicas}{Trees /\\logística}{NN corta\\+ calibrar}{Curvas\\lead-time}
  \vspace{0.35em}
  \begin{itemize}
    \item Unidad de análisis: cada \textbf{acercamiento completo} (no cada fila suelta).
    \item Métrica oficial del reto: mezcla F2 (casos graves) + error en high-risk.
    \item Oral: cuándo ayuda al operador y cuándo falla.
  \end{itemize}
\end{frame}

\begin{frame}{A · Métodos del curso y más allá}
  \MethodsSplit{Regresión / logística · Árboles / boosting · Validación cuidadosa · Ingeniería de variables}{Serie corta / NN · Calibración · Features orbitales · Opcional: incertidumbre / GP}
\end{frame}

\begin{frame}{A · Ángulos de originalidad (elegid uno)}
  \begin{enumerate}
    \item Curvas de \textbf{early-warning}: calidad vs días de antelación.
    \item \textbf{Calibración} e uncertainty (no solo score de competición).
    \item Transfer entre \texttt{mission\_id} (el paper documenta shift).
    \item Features físicas + ablación limpia.
    \item Clasificar high-risk + regresión condicional (como el diseño del score).
  \end{enumerate}
\end{frame}

\begin{frame}{A · Encaje Pedro / Mario}
  \FitDual{Máximo: espacio, ESA, física, anti-genérico. Lidera narrativa SSA y variables orbitales.}{Lidera métricas, rare events, calibración y comparación limpia de modelos.}{10}{7}
\end{frame}

\begin{frame}{A · Scorecard}
  \begin{columns}[T]
    \begin{column}{0.48\textwidth}
      \ScoreBar{Ambición / wow}{9}{ieblue}
      \ScoreBar{Datos / claridad}{9}{accent}
      \ScoreBar{Factibilidad 6 sem.}{8}{good}
      \ScoreBar{Afinidad Pedro}{10}{ieblue}
      \ScoreBar{Afinidad Mario}{7}{accent}
    \end{column}
    \begin{column}{0.48\textwidth}
      \ScoreBar{Oral}{9}{ieblue}
      \ScoreBar{Portfolio}{9}{accent}
      \ScoreBar{Alineación Molina}{9}{good}
      \ScoreBar{Diferenciación}{9}{ieblue}
      \ScoreBar{Riesgo acceso datos}{2}{good}
    \end{column}
  \end{columns}
  \vspace{0.3em}
  {\scriptsize Riesgo acceso: 1 = bajo (bueno), 10 = alto (malo). Aquí \textbf{2} = Zenodo ya disponible.}
\end{frame}

\begin{frame}{A · Pros y contras}
  \ProCon{Historia espacial muy fuerte\\Datos oficiales, descarga ya\\Factible en 6 semanas\\Buen oral físico / UQ\\Alineado con Molina}{Casos graves muy raros\\Competición 2019 ya cerrada\\Pocas ``fotos'' visuales\\Hay que aportar ángulo propio\\Oral puede ser muy técnico}
\end{frame}

\begin{frame}{A · Veredicto parcial}
  \Verdict{Muy fuerte como principal}{Ideal si la identidad es ``espacio / ESA / deeptech''. Reserva natural: PHBench.}
\end{frame}

% ============================================================
\section{Opción B --- PHBench}
% ============================================================

\begin{frame}{B · Portada --- ¿Qué startups llegan a Series A?}
  \OptionHeader{Finalista B}{Inversión / VC}{%
    Usamos señales públicas del día del lanzamiento en Product Hunt para estimar qué productos acaban levantando una ronda Series A.
  }
\end{frame}

\begin{frame}{B · Imagina esto (sin jerga)}
  \begin{itemize}
    \item Cada día salen cientos de apps y productos nuevos.
    \item Casi ninguno consigue una gran ronda (Series A) en $\sim$18 meses.
    \item Product Hunt es un ``escaparate'' público: votos, comentarios, equipo\ldots
    \item Un analista no puede mirarlo todo: necesita una \textbf{lista ordenada}.
    \item Nosotros entrenamos modelos que intentan hacer ese ranking con datos pasados.
  \end{itemize}
\end{frame}

\begin{frame}{B · Analogía cotidiana}
  \begin{center}
  \begin{tikzpicture}
    \node[box, text width=3.2cm] (a) {Muchas\\candidaturas};
    \node[box, text width=3.2cm, right=0.6cm of a] (b) {Pocas\\plazas};
    \node[boxdark, text width=3.2cm, right=0.6cm of b] (c) {Hay que\\priorizar};
    \draw[arr] (a)--(b)--(c);
  \end{tikzpicture}
  \end{center}
  \vspace{0.45em}
  Mismo tipo de problema que filtrar muchas solicitudes cuando solo unas pocas son ``éxito''.\\
  \textbf{No} es predecir si sube o baja una acción en bolsa.
\end{frame}

\begin{frame}{B · ¿Qué haría el machine learning?}
  \begin{itemize}
    \item Aprender patrones en lanzamientos antiguos (votos, temas, señales del día).
    \item Predecir probabilidad de Series A --- éxito \textbf{rarísimo} ($\sim$1 de cada 127).
    \item Evaluar como ranking: ``de los 50 que marcamos, ¿cuántos acertamos?''
    \item Cuidado crítico: no usar información del \textbf{futuro} (fugas temporales).
  \end{itemize}
\end{frame}

\begin{frame}{B · Glosario rápido}
  \begin{description}[font=\bfseries\scriptsize,style=nextline,leftmargin=1.5cm]
    \item[Series A] Primera gran ronda institucional de una startup.
    \item[Product Hunt] Comunidad donde se lanzan y votan productos digitales.
    \item[F0.5] Métrica primaria del benchmark: prioriza precisión (mejor pocos aciertos buenos).
    \item[Lift@k] Cuánto mejor que el azar es tu top-$k$.
    \item[Fuga] Variable que ``mira al futuro'' (p.\ ej.\ seguidores medidos años después).
  \end{description}
\end{frame}

\begin{frame}{B · Pregunta de investigación}
  \begin{block}{Principal}
    Con señales del día del launch, ¿qué modelos ordenan mejor los candidatos a Series A en 18 meses, respetando el orden temporal?
  \end{block}
  \begin{block}{Secundarias (elegid 1 como gancho)}
    \scriptsize
    (i) Ranking semanal tipo ``cola de analista''. \quad
    (ii) Ablación honesta de features con fuga. \quad
    (iii) Calibración + curvas precisión-capacidad. \quad
    (iv) ¿Cambian los patrones por ciclo de mercado?
  \end{block}
\end{frame}

\begin{frame}{B · Datos (hechos concretos)}
  \begin{itemize}
    \item $\sim$67\,292 launches; 528 positivos (0{,}78\%); tabular anonimizado; CC BY 4.0.
    \item \textbf{Sin} tagline/descripción cruda en el release público --- no pitch SBERT en week 1.
    \item Champion test: AP 0{,}037 ($\sim$4{,}7$\times$ random), F$_{0.5}$ 0{,}097 (val inflada).
    \item HF gated \textbf{o} email \texttt{benchmark@vela.partners} · paper arXiv:2605.02974.
  \end{itemize}
  \begin{block}{Riesgo operativo}
    Sin CSV mañana $\rightarrow$ no primary. Fallback si el gate falla: \textbf{VCBench} (no dual-core).
  \end{block}
\end{frame}

\begin{frame}{B · Métricas y trampas documentadas}
  \begin{columns}[T]
    \begin{column}{0.48\textwidth}
      \textbf{Métricas oficiales}
      \begin{itemize}
        \item Primaria: \textbf{F0.5}
        \item También: P@50/100/200, AP, lift
      \end{itemize}
    \end{column}
    \begin{column}{0.48\textwidth}
      \textbf{Trampas}
      \begin{itemize}
        \item Fuga en \texttt{maker\_followers}
        \item Cohort 2025 parcialmente resuelto
        \item Sesgo: PH $\neq$ universo startups
      \end{itemize}
    \end{column}
  \end{columns}
  \vspace{0.4em}
  Hablar de estas trampas en el oral = señal de madurez (no es un bug vuestro: está en el paper).
\end{frame}

\begin{frame}{B · Qué se ha intentado ya (benchmark 2026)}
  \begin{itemize}
    \item Benchmark fresco (Ihlamur et al., arXiv:2605.02974) con baselines públicos.
    \item \textbf{Champion:} ensemble de árboles (Top-3) --- test AP 0{,}037 / F$_{0.5}$ 0{,}097 ($\sim$4{,}7$\times$ random).
    \item Baseline oficial: regresión logística (FS4) --- test AP 0{,}024 / F$_{0.5}$ 0{,}045.
    \item Dominan: gradient boosting + interacciones + pesos de clase + calibración isotónica.
    \item LLMs zero-shot (Gemini, números anonimizados): a veces ganan en F$_{0.5}$ puntual, \textbf{pierden en AP}; modelo más grande $\neq$ mejor.
  \end{itemize}
\end{frame}

\begin{frame}{B · Qué queda abierto (ángulo propio)}
  \begin{itemize}
    \item Holdout \textbf{temporal} (train $\leq$2023 $\to$ eval 2024--25) --- lo dejan los autores como future work.
    \item Calibración hecha bien (ellos calibran en el mismo val que miden).
    \item Ablación honesta de fuga en followers; lift@$k$ con presupuesto de analista.
    \item Brazo tabular PyTorch vs árboles; opcional: LLM como ``picker'' + ML como ranker.
    \item \textbf{No} original: ``corrimos XGBoost'' ni ``los LLMs hacen VC'' sin protocolo.
  \end{itemize}
\end{frame}

\begin{frame}{B · Cómo lo haríamos (6 semanas)}
  \PipelineFive{Acceso +\\EDA}{Limpieza\\fugas}{LR / trees\\+ ranking}{Calibrar\\+ utilidades}{Errores +\\ciclos}
\end{frame}

\begin{frame}{B · Métodos del curso y más allá}
  \MethodsSplit{Logística regularizada · Boosting · Validación temporal · Métricas de ranking}{Calibración · Ablaciones de fuga · Ensambles · Análisis por cohorte / ciclo}
\end{frame}

\begin{frame}{B · Ángulos de originalidad (elegid uno)}
  \begin{enumerate}
    \item Framing \textbf{rank-the-weekly-queue} + lift@$k$ / presupuesto semanal.
    \item Ablación honesta de features con fuga temporal.
    \item Holdout temporal (train $\leq$2023, eval 2024--25).
    \item Calibración anidada + curvas precisión-capacidad.
    \item Topics / stats de superficie (no embeddings de texto crudo).
  \end{enumerate}
\end{frame}

\begin{frame}{B · Encaje Pedro / Mario}
  \FitDual{VC innovador (no bolsa genérica). Portfolio fuerte. Menos ``física'', más narrativa de mercado.}{Máximo: finanzas + ranking. Datos claros, comparación limpia, métricas de negocio.}{8}{10}
\end{frame}

\begin{frame}{B · Scorecard}
  \begin{columns}[T]
    \begin{column}{0.48\textwidth}
      \ScoreBar{Ambición / wow}{9}{ieblue}
      \ScoreBar{Datos / claridad}{8}{accent}
      \ScoreBar{Factibilidad 6 sem.}{7}{accent}
      \ScoreBar{Afinidad Pedro}{8}{accent}
      \ScoreBar{Afinidad Mario}{10}{good}
    \end{column}
    \begin{column}{0.48\textwidth}
      \ScoreBar{Oral}{8}{ieblue}
      \ScoreBar{Portfolio}{10}{good}
      \ScoreBar{Alineación Molina}{6}{accent}
      \ScoreBar{Diferenciación}{9}{ieblue}
      \ScoreBar{Riesgo acceso datos}{7}{bad}
    \end{column}
  \end{columns}
  \vspace{0.3em}
  {\scriptsize Riesgo acceso 7/10: HF gated + email. Mitigable si confirmáis acceso esta noche.}
\end{frame}

\begin{frame}{B · Pros y contras}
  \ProCon{Poco genérico / fresco 2026\\Encaja muy bien con Mario\\Baselines y leaderboard listos\\Historia fuerte en entrevista\\Métricas de ranking naturales}{Acceso a datos por confirmar\\Desbalance brutal (0{,}78\%)\\Fácil contaminar con futuro\\Menos deeptech físico\\Ética / sesgo en el oral}
\end{frame}

\begin{frame}{B · Veredicto parcial}
  \Verdict{Muy fuerte si hay datos ya}{Ideal si la identidad es ``VC / portfolio / ranking''. Sin acceso confirmado, mejor como reserva.}
\end{frame}

% ============================================================
\section{Cara a cara}
% ============================================================

\begin{frame}{Cara a cara · Identidad del proyecto}
  \noindent
  \begin{minipage}[t]{0.48\textwidth}
    \begin{beamercolorbox}[wd=\textwidth,rounded=true,sep=8pt,center]{block title}
      \textbf{ESA}\\[0.35em]
      {\scriptsize\normalfont ``Trabajamos con datos de la Agencia Espacial Europea para avisar antes del riesgo.''}
    \end{beamercolorbox}
  \end{minipage}\hfill
  \begin{minipage}[t]{0.48\textwidth}
    \begin{beamercolorbox}[wd=\textwidth,rounded=true,sep=8pt,center]{block title}
      \textbf{PHBench}\\[0.35em]
      {\scriptsize\normalfont ``Ordenamos la cola semanal de launches para no perder tiempo en early-stage.''}
    \end{beamercolorbox}
  \end{minipage}
  \vspace{0.55em}
  Pregunta guía: \textbf{¿qué frase queréis decir en la primera frase del oral?}
\end{frame}

\begin{frame}{Cara a cara · Tabla resumen (1--10)}
  \tiny
  \renewcommand{\arraystretch}{1.2}
  \begin{tabular}{@{}lcc@{}}
    \toprule
    Criterio & ESA & PHBench \\
    \midrule
    Ambición / wow & 9 & 9 \\
    Calidad / claridad datos & 9 & 8 \\
    Factibilidad 6 semanas & 8 & 7 \\
    Afinidad Pedro & 10 & 8 \\
    Afinidad Mario & 7 & 10 \\
    Articulable con temario & 9 & 9 \\
    Defendible en oral & 9 & 8 \\
    Portfolio laboral & 9 & 10 \\
    Alineación Molina (física / UQ) & 9 & 7 \\
    Diferenciación vs genérico & 9 & 8 \\
    Riesgo de verse genérico ($\uparrow$=peor) & \textbf{2} & 4 \\
    Riesgo de no poder empezar mañana ($\uparrow$=peor) & \textbf{2} & \textbf{7} \\
    \bottomrule
  \end{tabular}
\end{frame}

\begin{frame}{Cara a cara · Quién gana cada eje}
  \begin{columns}[T]
    \begin{column}{0.48\textwidth}
      \begin{block}{Gana ESA}
        \scriptsize
        Acceso inmediato · Pedro/espacio · Molina/física/calibración · narrativa ESA · menor fricción esta noche
      \end{block}
    \end{column}
    \begin{column}{0.48\textwidth}
      \begin{block}{Gana PHBench}
        \scriptsize
        Mario/ranking · portfolio VC · frescura 2026 · métricas de negocio · entrevista ``data for investing''
      \end{block}
    \end{column}
  \end{columns}
  \vspace{0.45em}
  Empate real en ambición y diferenciación. El desempate práctico es: \textbf{¿datos en disco mañana?}
\end{frame}

\begin{frame}{Cara a cara · Tipo de ML}
  \begin{center}
  \begin{tikzpicture}
    \node[box, text width=5.0cm] (a) {\textbf{ESA}\\serie corta por evento\\regresión + rare events\\física / calibración};
    \node[box, text width=5.0cm, right=0.55cm of a] (b) {\textbf{PHBench}\\tabular + ranking\\clasificación extrema\\validación temporal / fugas};
    \draw[arr] (a)--(b);
  \end{tikzpicture}
  \end{center}
  \vspace{0.45em}
  Ambas usan el temario (trees, logística, validación). El ``extra'' es distinto: física/UQ vs ranking/ética de datos.
\end{frame}

\begin{frame}{Cara a cara · Oral (50\%)}
  \begin{columns}[T]
    \begin{column}{0.48\textwidth}
      \textbf{ESA --- preguntas típicas}
      \begin{itemize}
        \item ¿Qué es un CDM?
        \item ¿Por qué no basta el accuracy?
        \item ¿Qué es LRP y por qué es fuerte?
        \item ¿Calibración vs score?
      \end{itemize}
    \end{column}
    \begin{column}{0.48\textwidth}
      \textbf{PHBench --- preguntas típicas}
      \begin{itemize}
        \item ¿Causalidad o correlación?
        \item ¿Sesgo de Product Hunt?
        \item ¿Qué fugas evitasteis?
        \item ¿Por qué F0.5 y no F1?
      \end{itemize}
    \end{column}
  \end{columns}
\end{frame}

\begin{frame}{Cara a cara · Proposal de 2 páginas}
  \begin{itemize}
    \item \textbf{ESA:} más fácil de escribir ya (links Zenodo + pregunta + baseline LRP + ángulo early-warning).
    \item \textbf{PHBench:} igual de redactable, pero debéis mencionar plan de acceso y métricas F0.5 / P@k.
    \item En ambos casos: un solo gancho de novedad, no tres.
  \end{itemize}
\end{frame}

\begin{frame}{Cara a cara · Riesgos que matan el proyecto}
  \ProCon{\textbf{ESA}\\Quedarse en ``rehacer la competición''\\sin ángulo propio\\$\rightarrow$ oral flojo de originalidad}{\textbf{PHBench}\\Empezar sin CSV\\o ignorar fugas temporales\\$\rightarrow$ resultados no creíbles}
\end{frame}

\begin{frame}{Cara a cara · Visual de trade-offs}
  \begin{center}
  \begin{tikzpicture}[scale=0.95]
    % axes
    \draw[->, thick] (0,0) -- (8.2,0) node[right, font=\scriptsize] {Afinidad Mario};
    \draw[->, thick] (0,0) -- (0,5.2) node[above, font=\scriptsize] {Afinidad Pedro};
    \node[boxdark, fill=esa, text width=2.4cm] at (5.0,4.3) {ESA\\{\scriptsize Pedro 10 · Mario 7}};
    \node[boxdark, fill=phb, text width=2.4cm] at (7.0,3.4) {PHBench\\{\scriptsize Pedro 8 · Mario 10}};
    \node[font=\tiny, gray] at (4,-0.45) {ambas son ambiciosas; el eje es identidad};
  \end{tikzpicture}
  \end{center}
\end{frame}

% ============================================================
\section{Decisión}
% ============================================================

\begin{frame}{Escenarios --- elegid el vuestro}
  \tiny
  \renewcommand{\arraystretch}{1.25}
  \begin{tabular}{@{}p{7.2cm}l@{}}
    \toprule
    Situación & Elegir \\
    \midrule
    Queréis identidad ``espacio / ESA / deeptech'' & \textbf{ESA} \\
    Queréis identidad ``VC / portfolio / ranking'' & \textbf{PHBench}* \\
    Proposal hoy y PHBench aún sin acceso & \textbf{ESA} + PHBench reserva \\
    Mario debe brillar en métricas de negocio & \textbf{PHBench}* \\
    Maximizar alineación Molina (física / UQ) & \textbf{ESA} \\
    Teméis parecer ``otra competición Kelvins'' & PHBench* o ESA con ángulo claro \\
    \bottomrule
  \end{tabular}\\[0.4em]
  {\scriptsize *Solo si el acceso a datos está confirmado (HF o email).}
\end{frame}

\begin{frame}{Recomendación opinada (research)}
  \Verdict{Default: ESA primary · PHBench reserva}{Confianza $\sim$72\%. Con peso oral (50\%) + ambición (25\%), ESA gana por margen estrecho pero real.}
  \vspace{0.25em}
  \begin{block}{Por qué}
    Dato listo · física/UQ · Pedro a tope. Mario no pierde: lidera rare events, calibración y ``cuándo el modelo miente''.
  \end{block}
  \begin{block}{Vista disidente ($\sim$28\%)}
    Si el portfolio VC importa más esta noche \textbf{y} hay acceso a datos $\rightarrow$ PHBench primary.
  \end{block}
\end{frame}

\begin{frame}{Híbridos a evitar (falsos compromisos)}
  \begin{itemize}
    \item ESA + PHBench en \textbf{un} proyecto $\rightarrow$ dos RQs, oral diluido.
    \item Metáfora ``Series A $\approx$ riesgo orbital'' $\rightarrow$ deck, no ciencia.
    \item Mitad-mitad ``por equidad'' $\rightarrow$ equidad de \textbf{roles}, no de tema.
    \item Transformer ``porque SOTA'' como núcleo $\rightarrow$ ambición falsa.
  \end{itemize}
  \vspace{0.35em}
  \begin{block}{Compromiso sano}
    Un primary + un backup en 3 líneas + roles asimétricos claros.
  \end{block}
\end{frame}

\begin{frame}{Checklist proposal (gates duros)}
  \begin{columns}[T]
    \begin{column}{0.48\textwidth}
      \begin{itemize}
        \item Primary elegido (A o B)
        \item Backup en 3 líneas
        \item RQ falsable (1 frase)
        \item URL datos + estado de acceso
        \item Unidad de análisis clara
      \end{itemize}
    \end{column}
    \begin{column}{0.48\textwidth}
      \begin{itemize}
        \item Métrica rare-events (no accuracy)
        \item Split sin filtrar futuro
        \item Un solo ángulo propio
        \item Ancla temario sklearn + PyTorch
        \item Roles Pedro / Mario / oral
      \end{itemize}
    \end{column}
  \end{columns}
\end{frame}

\begin{frame}{Reglas de decisión (recordatorio)}
  \begin{itemize}
    \item Pedro puede vetar lo \textbf{genérico}.
    \item Mario puede vetar lo \textbf{mal definido}.
    \item Ambición alta, pero ejecutable en 6 semanas.
    \item Solo \textbf{un} gancho de novedad.
    \item Oral: Mario abre con datos/métricas; Pedro cierra con el porqué.
  \end{itemize}
\end{frame}

\begin{frame}{Scorecard para rellenar juntos}
  \tiny
  \renewcommand{\arraystretch}{1.35}
  \begin{tabular}{@{}clcccc@{}}
    \toprule
    & Opción & Pedro (1--5) & Mario (1--5) & Factib. & ¿Principal? \\
    \midrule
    A & ESA Collision Avoidance & & & & \\
    B & PHBench Series A & & & & \\
    \bottomrule
  \end{tabular}
  \vspace{0.55em}
  \normalsize
  Descartad como principal cualquier opción con factibilidad $<$ 3 \textbf{o} sin datos esta semana.
\end{frame}

\begin{frame}{Spike de 48 horas (si empatáis)}
  \noindent
  \begin{minipage}[t]{0.48\textwidth}
    \begin{beamercolorbox}[wd=\textwidth,rounded=true,sep=6pt,center]{block title}
      \scriptsize ESA: Zenodo + 2 plots + baseline LRP
    \end{beamercolorbox}
  \end{minipage}\hfill
  \begin{minipage}[t]{0.48\textwidth}
    \begin{beamercolorbox}[wd=\textwidth,rounded=true,sep=6pt,center]{block title}
      \scriptsize PHBench: pedir acceso HF + email hoy
    \end{beamercolorbox}
  \end{minipage}
  \vspace{0.45em}
  \begin{block}{Regla de desempate}
    Si en 24h no hay CSV de PHBench $\rightarrow$ ESA primary. Decidir con evidencia, no con vibes.
  \end{block}
\end{frame}

\begin{frame}{Plan inmediato tras elegir}
  \PipelineFive{Elegir\\principal}{Fijar\\reserva}{Bajar\\datos}{2--3\\gráficos}{Escribir\\proposal}
  \vspace{0.45em}
  \begin{block}{La proposal debe incluir}
    Pregunta precisa · motivación · dataset · exploración preliminar · método inicial · plan de 6 semanas.
  \end{block}
\end{frame}

\begin{frame}{Fuentes clave}
  \begin{columns}[T]
    \begin{column}{0.48\textwidth}
      \textbf{ESA}
      \begin{itemize}
        \item zenodo.org/records/4463683
        \item kelvins.esa.int/...
        \item arXiv:2008.03069
      \end{itemize}
    \end{column}
    \begin{column}{0.48\textwidth}
      \textbf{PHBench}
      \begin{itemize}
        \item phbench.com
        \item github.com/ihlamury/phbench
        \item arXiv:2605.02974
        \item HF: ihlamury/phbench
      \end{itemize}
    \end{column}
  \end{columns}
  \vspace{0.35em}
  {\scriptsize
  Comparativa: \texttt{research-brief-esa-vs-phbench.md}\\
  Deep A/B: \texttt{research-brief-esa-collision-avoidance-deep.md} · \texttt{research-brief-phbench-deep.md}}
\end{frame}

\begin{frame}{Cierre}
  \begin{center}
    {\Large\textbf{¿ESA o PHBench?}}\\[0.6em]
    Una principal. Una reserva.\\[0.4em]
    Luego: datos en disco $\rightarrow$ proposal.
  \end{center}
\end{frame}

\end{document}
